\documentclass[11pt]{article}
\usepackage[margin=1in]{geometry}
\usepackage{amsmath,amssymb,amsthm}
\usepackage[hidelinks]{hyperref}
\newtheorem{theorem}{Theorem}
\DeclareMathOperator{\Var}{Var}
\title{Failure of the discrete maximal-variation conjecture\\
\large MathDB 348574}
\author{Principia Math in Progress}
\date{12 September 2026}
\begin{document}
\maketitle

\paragraph{Task 1. Resolve the original conjecture.}
For a function on the integers, use the centered maximal function and
adjacent-difference variation
\[
Mf(n)=\sup_{r\in\mathbb Z_{\geq0}}
\frac{\sum_{k=-r}^r|f(n+k)|}{2r+1},
\qquad
\Var_p(f)^p=\sum_{n\in\mathbb Z}|f(n+1)-f(n)|^p.
\]
Conjecture~1, equation~(5.4), of Gonz\'alez-Riquelme and
Madrid~\cite{source} proposes, for \(1/2<p\leq1\) and
\(f\in\ell^p(\mathbb Z)\),
\begin{equation}\label{eq:conjecture}
\Var_p(Mf)\leq C_p\Var_p(f),\qquad
C_p^p=\sum_{j=0}^{\infty}
\left(\frac{2}{(2j+1)(2j+3)}\right)^p.
\end{equation}
This is false for every \(1/2<p<1\). We first show that no finite
replacement constant works in that range, then give one fixed
three-point example refuting the specific proposed constant for all
those exponents.

\begin{theorem}\label{thm:unbounded}
For each \(1/2<p<1\),
\[
\sup_{\substack{f\ne0\\ f\text{ finitely supported}}}
\frac{\Var_p(Mf)}{\Var_p(f)}=\infty.
\]
\end{theorem}
\begin{proof}
Let \(N\geq1\) be an integer and put
\(f_N=\mathbf1_{[-N,N]\cap\mathbb Z}\).
There are exactly two unit jumps, so \(\Var_p(f_N)^p=2\).
We claim that
\begin{equation}\label{eq:block}
Mf_N(q)=
\begin{cases}
1,&|q|\leq N,\\[1mm]
(2N+1)/(2|q|+2N+1),&|q|>N.
\end{cases}
\end{equation}
For \(q>N\), the first radius capturing \(j\) support points is
\(q-N+j-1\), with average
\[
\frac{j}{2q-2N+2j-1},\qquad 1\leq j\leq2N+1.
\]
These fractions strictly increase with \(j\), because \(2q-2N-1>0\).
Once the whole support is captured, every larger radius decreases the
average. This proves the claim; symmetry and radius zero handle the
other centers.

For \(N+1\leq q\leq2N\), formula~\eqref{eq:block} gives
\[
Mf_N(q)-Mf_N(q+1)
=\frac{2(2N+1)}{(2q+2N+1)(2q+2N+3)}
\geq\frac4{63N}.
\]
Indeed, the numerator is at least \(4N\), and the denominator is at
most \((6N+1)(6N+3)\leq63N^2\).
There are \(N\) such jumps on each side. Consequently
\[
\left(\frac{\Var_p(Mf_N)}{\Var_p(f_N)}\right)^p
\geq N\left(\frac4{63N}\right)^p
=\left(\frac4{63}\right)^pN^{1-p}\longrightarrow\infty.
\]
Each \(f_N\) belongs to \(\ell^p\), and each output variation is finite:
for fixed \(N\), the tail jumps are \(O(q^{-2})\) and \(2p>1\).
\end{proof}

\paragraph{A fixed counterexample to the proposed constant.}
Take \(f=f_1=\mathbf1_{\{-1,0,1\}}\). Its maximal function is
\[
Mf(n)=1\quad(|n|\leq1),\qquad
Mf(n)=\frac3{2|n|+3}\quad(|n|\geq2).
\]
The jump sizes on either side are
\[
a_0=\frac47,\qquad
a_j=\frac6{(2j+5)(2j+7)}\quad(j\geq1).
\]
Set also
\[
b_j=\frac2{(2j+1)(2j+3)}\quad(j\geq0).
\]
Both \(a\) and \(b\) are decreasing positive sequences of total mass
one. For \(m\geq1\), their partial sums are
\[
A_m=\sum_{j<m}a_j=1-\frac3{2m+5},\qquad
B_m=\sum_{j<m}b_j=1-\frac1{2m+1},
\]
so
\begin{equation}\label{eq:major}
B_m-A_m=\frac{2(2m-1)}{(2m+1)(2m+5)}>0.
\end{equation}
We make the infinite-sequence concavity step explicit.
For any decreasing probability sequence \(z=(z_j)\), put
\[
K_z(t)=\sum_j\min(z_j,t)\qquad(t>0).
\]
If \(Z_m=\sum_{j<m}z_j\) and \(Z_0=0\), then
\[
\sum_j(z_j-t)_+=\max_{m\geq0}(Z_m-mt),\qquad
K_z(t)=1-\sum_j(z_j-t)_+.
\]
The maximum is attained at the finite number of entries exceeding \(t\).
Equation~\eqref{eq:major} therefore implies \(K_a(t)\geq K_b(t)\).
This inequality is strict for \(4/7<t<2/3\): every \(a_j<t\), so
\(K_a(t)=1\), while \(b_0=2/3>t\) gives \(K_b(t)<1\).

For \(0<p<1\), direct integration gives
\[
z^p=p(1-p)\int_0^\infty \min(z,t)t^{p-2}\,dt.
\]
Tonelli's theorem and the preceding strict interval now show
\[
\sum_j a_j^p>\sum_j b_j^p\qquad(1/2<p<1).
\]
Both sums are finite because their terms are \(O(j^{-2p})\).
Finally,
\[
\left(\frac{\Var_p(Mf)}{\Var_p(f)}\right)^p
=\frac{2\sum_ja_j^p}{2}
>\sum_jb_j^p=C_p^p.
\]
Thus the same three-point function refutes~\eqref{eq:conjecture}
for every exponent in the open interval.

\paragraph{Scope.}
At \(p=1\) both jump sums equal one, so these examples give equality.
The proof refutes the universally quantified conjecture and all finite
variation-to-variation constants for \(1/2<p<1\); it does not resolve
the general endpoint \(p=1\).
The MathDB strings \(p\ne(1/2,1]\) and the malformed \(\ell^p\)
are transcription defects: the source and the record's explanatory
sentence fix the intended range and space.
The source's proved \(\ell^p\)-to-variation bound is a different estimate
and is not contradicted.

\paragraph{Task 2. Quantify the dependence on support size.}
Determine the sharp growth of the best variation-to-variation constant
when the input is supported on an interval of at most \(L\) lattice
points, for fixed \(1/2<p<1\).

\paragraph{Task 3. Resolve the endpoint variation bound.}
Prove or refute \(\Var_1(Mf)\leq\Var_1(f)\) for general
\(f\in\ell^1(\mathbb Z)\), extending beyond the already known
indicator-function case.

\paragraph{Task 4. Study radius truncation.}
For the maximal operator restricted to radii \(r\leq R\), determine
the optimal variation constant and its growth as \(R\to\infty\)
when \(1/2<p<1\).

\paragraph{Task 5. Find a replacement estimate.}
Determine sharp bounds combining input variation with a measure of
support size, mass or decay that controls the long-block obstruction
while retaining the established \(\ell^p\)-to-variation inequality.

\begin{thebibliography}{9}
\bibitem{source}
C. Gonz\'alez-Riquelme and J. Madrid,
\emph{Sharp inequalities for maximal operators on finite graphs, II},
J. Math. Anal. Appl. \textbf{506} (2022), 125647,
\href{https://doi.org/10.1016/j.jmaa.2021.125647}{doi:10.1016/j.jmaa.2021.125647};
\href{https://arxiv.org/abs/2011.02630}{arXiv:2011.02630}, Conjecture~1.
\end{thebibliography}
\end{document}
